% !TeX root = ../../Skript.tex
\cohead{\Large\textbf{Betrag eines Vektors}}
\fakesubsection{Betrag eines Vektors}

\begin{minipage}{\textwidth}
	\adjustbox{valign=t, margin=0ex 0ex 5ex 0ex}{\begin{minipage}{0.55\textwidth-5ex}\raggedright
			Ein Vektor hat eine Richtung und eine Länge. Die Länge bezeichnet man auch als den Betrag des Vektors. Das mathematische Symbol für den Betrag sind zwei vertikale Striche: \(\abs{\vv{v}}\)
			
			Die Länge eines Vektors entspricht der Raumdiagonalen eines Quaders mit den Seitenlängen, die den Beträgen der Komponenten des Vektors entsprechen. 
			
			\textcolor{loes}{Berechnen der Diagonalen \(d\) des Rechtecks mit Seitenlängen \(x_1\) und \(x_2\):}
			
			\(\textcolor{loes}{d^2=x_1^2+x_2^2}\)
			
			\textcolor{loes}{Damit ergibt sich für den Betrag:}
			
			\begin{align*}
				\textcolor{loes}{\abs{\vv{v}}\,}&\textcolor{loes}{=d^2+x_3^2}\\
				&\textcolor{loes}{=x_1^2+x_2^2+x_3^2}
			\end{align*}
			
			\bigskip
			
			Der Betrag ist eine Abbildung, die aus dem \(\R^3\) oder \(\R^3\) auf \(\R\) abbildet: \(\abs{\cdot}:\R^3\rightarrow\R\)
	\end{minipage}}%
	\adjustbox{valign=t}{\begin{minipage}{0.45\textwidth}
			\includegraphics[width=\textwidth]{\vektor/pics/Betrag.png}
	\end{minipage}}%
\end{minipage}

\begin{tcolorbox}
	\textbf{Betrag eines Vektors}
	
	\textcolor{loestc}{Der Betrag bzw. die Länge eines Vektors lässt sich wie folgt berechnen:}
	
	\[\textcolor{loestc}{\abs{\vv{v}}=\abs{\spvvv{x_1}{x_2}{x_3}}=\sqrt{x_1^2+x_2^2+x_3^2}\text{ im }\R^3\text{ und }\abs{\vv{v}}=\abs{\spvv{x_1}{x_2}}=\sqrt{x_1^2+x_2^2}\text{ im }\R^2}\]
	
	\bigskip
	
\end{tcolorbox}

\begin{tcolorbox}
	\textbf{Einheitsvektoren}
	
	\textcolor{loestc}{Der Einheitsvektor eines Vektors \(\vv{v}\) für \(\vv{v}\neq \vv{o}\) ist der Vektor, der die gleiche Richtung hat und die Länge 1:}
	
	\[\textcolor{loestc}{\vv{v_0}=\frac{\vv{v}}{\abs{\vv{v}}}}\]
	
	\bigskip
	
\end{tcolorbox}
Hinweis: Man spricht auch vom Normieren eines Vektors, wenn man seine Länge auf 1 bringt. Es gibt verschiedene Notationen für einen Einheitsvektor, z.B. \(\vv{v_0}=\vv{e_v}=\hat{v}\)

Die Einheitsvektoren parallel zu den Achsen bezeichnet man mit einem \(e\), z.B. \(\vv{e_1}=\spvvv{1}{0}{0}\)
%\begin{Exercise}[title={\raggedright xxxxxxxxxxxxxxxxxxxxxxxxxxxxxxxxxxxxxxxxxxxxxxxxxxxx}, label=VektorenBetragA1]
%	
%\end{Exercise}
%\newpage
%%%%%%%%%%%%%%%%%%%%%%%%%%%%%%%%%%%%%%%%%%
%\begin{Answer}[ref=VektorenBetragA1]
%	%TODO
%\end{Answer}